% ============================================================
%  FreshStock
%  Universidad de San Buenaventura — Ingeniería Multimedia
% ============================================================

\documentclass[journal]{IEEEtran}


\usepackage[utf8]{inputenc}
\usepackage[T1]{fontenc}
\usepackage[spanish,es-tabla]{babel}
\usepackage{graphicx}
\usepackage[hidelinks]{hyperref}
\usepackage{tabularx}
\usepackage{booktabs}
\usepackage{placeins}
\usepackage{microtype}
\usepackage{xcolor}
\usepackage{array}
\usepackage{enumitem}
\usepackage{cite}
\usepackage{balance}   


%  CONFIGURACION DE HYPERREF

\hypersetup{
	colorlinks        = true,
	linkcolor         = black,
	citecolor         = black,
	urlcolor          = blue!70!black,
	pdftitle          = {FreshStock: Sistema de Gestion de Inventarios},
	pdfauthor         = {Zambrano, Chavez, Acosta},
	pdfsubject        = {Ingenieria de Software -- PWA},
	pdfkeywords       = {inventario, caducidad, PWA, PEPS, base de datos},
	bookmarksnumbered = true,
	breaklinks        = true
}


%  COMANDO PARA IMAGEN CON MARCADOR DE POSICION

\newcommand{\figura}[4]{%
	\IfFileExists{#3}{%
		\includegraphics[width=#1,height=#2,keepaspectratio]{#3}%
	}{%
		\fboxsep=0pt
		\colorbox{gray!18}{%
			\parbox[c][#2][c]{#1}{%
				\centering\color{gray!60}\small [Imagen: #4]%
			}%
		}%
	}%
}


%  SEPARACION SILABICA

\hyphenation{op-tical net-works semi-con-duc-tor Fresh-Stock
	ca-du-ci-dad in-ven-ta-rio tra-za-bi-li-dad}

%  MACROS

\newcommand{\app}{\textit{FreshStock}}
\newcommand{\peps}{\textbf{PEPS}}
\newcommand{\pwa}{\textit{PWA}}


%  INICIO DEL DOCUMENTO

\begin{document}
	
	% -----------------------------------------------------------
	%  CABECERA
	% -----------------------------------------------------------
	\title{%
		FreshStock: Sistema de Gesti\'{o}n de Inventarios\\
		y Control de Caducidad para Peque\~{n}os Comercios de Alimentos%
	}
	
	\author{%
		Brian~Jesid~Zambrano~Poveda,\;
		Juan~Pablo~Ch\'{a}ves~Estupi\~{n}\'{a}n,\;
		Sadid~Enrique~Acosta~Osorio%

	}
	
	\markboth{%
		Universidad de San Buenaventura --- Ingenier\'{i}a Multimedia, 2026%
	}{%
		Zambrano \MakeLowercase{\textit{et al.}}: \app{} ---
		Gesti\'{o}n de Inventarios y Caducidad%
	}
	
	\maketitle
	
	% -----------------------------------------------------------
	%  RESUMEN
	% -----------------------------------------------------------
	\begin{abstract}
		Este documento describe el dise\~{n}o y la planificaci\'{o}n de \app{},
		una Aplicaci\'{o}n Web Progresiva (\pwa{}) orientada a la gesti\'{o}n de
		inventario y al control de fechas de caducidad para peque\~{n}os y
		medianos negocios minoristas de alimentos. La problem\'{a}tica central
		abordada son las p\'{e}rdidas econ\'{o}micas y las ineficiencias operativas
		ocasionadas por la ausencia de trazabilidad en tiempo real sobre
		productos perecederos. La soluci\'{o}n propuesta integra el m\'{e}todo
		\peps{} (Primero en Entrar, Primero en Salir), alertas autom\'{a}ticas
		de caducidad y un sistema de auditor\'{i}a inmutable. El documento cubre
		el an\'{a}lisis del problema, el estado del arte, una propuesta funcional
		detallada, el dise\~{n}o del flujo de navegaci\'{o}n, el desarrollo de un
		prototipo interactivo de alta fidelidad y el modelado de la base de
		datos relacional que soporta la operativa del sistema.
	\end{abstract}
	
	\begin{IEEEkeywords}
		Gesti\'{o}n de inventario, control de fechas de caducidad, \pwa{},
		m\'{e}todo \peps{}, modelo Entidad-Relaci\'{o}n, base de datos, trazabilidad.
	\end{IEEEkeywords}
	

	%  SECCION 1 — PLANTEAMIENTO DEL PROBLEMA

	\section{Planteamiento del Problema}
	\label{sec:problema}
	
	\IEEEPARstart{E}{n} el sector de los peque\~{n}os y medianos comercios de
	alimentos ---minimarkets, panader\'{i}as, fruter\'{i}as y similares--- el
	control de inventario es un factor determinante para la rentabilidad
	del negocio. A diferencia de las grandes cadenas, que cuentan con
	sistemas ERP especializados, estos establecimientos dependen con
	frecuencia de m\'{e}todos manuales o de hojas de c\'{a}lculo descentralizadas,
	lo que introduce riesgos operativos considerables.
	
	% ------- 1.1 -----------------------------------------------
	\subsection{Situaci\'{o}n Actual y Problem\'{a}tica}
	\label{subsec:situacion}
	
	La problem\'{a}tica central radica en la \textbf{falta de trazabilidad
		y control en tiempo real sobre los productos perecederos}.
	Las principales dificultades identificadas son:
	
	\begin{itemize}[leftmargin=*, label=\textbullet, itemsep=2pt]
		\item \textbf{Desconocimiento del stock real:} la disponibilidad
		exacta solo se confirma mediante verificaci\'{o}n f\'{i}sica,
		generando interrupciones en el servicio.
		
		\item \textbf{Gesti\'{o}n ineficiente de caducidades:} los productos
		se organizan por orden de llegada, permitiendo que los
		art\'{i}culos m\'{a}s antiguos permanezcan al fondo y caduquen
		sin ser rotados oportunamente.
		
		\item \textbf{Reabastecimiento reactivo:} las compras se realizan
		\'{u}nicamente cuando la escasez se detecta visualmente,
		imposibilitando la planificaci\'{o}n anticipada.
	\end{itemize}
	
	% ------- 1.2 -----------------------------------------------
	\subsection{Perfiles de Usuario Objetivo}
	\label{subsec:usuarios}
	
	El sistema est\'{a} dirigido a tres perfiles operativos diferenciados:
	
	\begin{itemize}[leftmargin=*, label=\textbullet, itemsep=2pt]
		\item \textbf{Administradores:} requieren visibilidad global,
		reportes de merma y control de usuarios del sistema.
		
		\item \textbf{Encargados de almac\'{e}n:} responsables del registro
		de entradas y de la gesti\'{o}n de lotes por fecha de caducidad.
		
		\item \textbf{Vendedores:} necesitan registrar salidas y consultar
		disponibilidad de forma \'{a}gil durante la atenci\'{o}n al cliente.
	\end{itemize}
	
	% ------- 1.3 -----------------------------------------------
	\subsection{Consecuencias de No Resolver el Problema}
	\label{subsec:consecuencias}
	
	La gesti\'{o}n actual genera impactos directos sobre la viabilidad del
	negocio, resumidos en la Tabla~\ref{tab:impacto}.
	
	\begin{table}[htbp]
		\renewcommand{\arraystretch}{1.35}
		\caption{Impacto negativo de la gesti\'{o}n manual de inventarios}
		\label{tab:impacto}
		\centering
		\begin{tabular}{>{\bfseries}p{2.4cm} p{5.2cm}}
			\toprule
			Impacto & Descripci\'{o}n \\
			\midrule
			P\'{e}rdidas econ\'{o}micas
			& Desechar productos por caducidad representa
			una p\'{e}rdida directa e irrecuperable de capital. \\[3pt]
			Da\~{n}o reputacional
			& Vender productos pr\'{o}ximos a vencer genera
			inconformidad en el cliente y deteriora la imagen
			del negocio. \\[3pt]
			Ineficiencia operativa
			& El personal invierte tiempo en conteos manuales
			en lugar de dedicarlo a la atenci\'{o}n al cliente. \\
			\bottomrule
		\end{tabular}
	\end{table}
	

	%  SECCION 2 — JUSTIFICACION DE LA SOLUCION

	\section{Justificaci\'{o}n de la Soluci\'{o}n}
	\label{sec:justificacion}
	
	% ------- 2.1 -----------------------------------------------
	\subsection{El Impacto del Desperdicio y la Necesidad de Control}
	\label{subsec:impacto}
	
	A nivel global, la ineficiencia en la cadena de suministro y la
	mala gesti\'{o}n de inventarios de perecederos contribuyen
	significativamente a las p\'{e}rdidas econ\'{o}micas y al desperdicio de
	alimentos~\cite{FAO:2019}. En el contexto de los peque\~{n}os comercios,
	implementar pr\'{a}cticas estandarizadas de gesti\'{o}n de inventario
	constituye un factor cr\'{i}tico de \'{e}xito que con frecuencia se descuida
	por limitaciones t\'{e}cnicas y presupuestales~\cite{PMI:2019}.
	
	\app{} se propone cerrar esta brecha mediante la democratizaci\'{o}n
	del acceso a herramientas de control avanzado. Se estima que la
	implementaci\'{o}n puede ayudar a que no haya más desperdicios
	por caducidad, transformando un modelo \emph{reactivo} en uno
	\emph{proactivo}.
	
	% ------- 2.2 -----------------------------------------------
	\subsection{Viabilidad Tecnol\'{o}gica: PWA}
	\label{subsec:pwa}
	
	Se propone el desarrollo de una \textbf{Aplicaci\'{o}n Web Progresiva
		(\pwa{})} como arquitectura id\'{o}nea, por dos razones fundamentales:
	
	\begin{itemize}[leftmargin=*, label=\textbullet, itemsep=2pt]
		\item \textbf{Despliegue de bajo costo:} opera desde los
		\textit{smartphones} o computadores ya disponibles en
		el negocio, sin hardware especializado.
		
		\item \textbf{Sincronizaci\'{o}n y concurrencia:} m\'{u}ltiples usuarios
		registran entradas y salidas simult\'{a}neamente, manteniendo
		una \'{u}nica fuente de verdad centralizada.
	\end{itemize}
	
	% ------- 2.3 -----------------------------------------------
	\subsection{Matriz de Operaciones CRUD}
	\label{subsec:crud}
	
	Para garantizar la inmutabilidad de los historiales, el sistema
	restringe eliminaci\'{o}n y actualizaci\'{o}n en entidades sensibles,
	tal como muestra la Tabla~\ref{tab:crud}.
	
	\begin{table}[htbp]
		\renewcommand{\arraystretch}{1.3}
		\caption{Operaciones CRUD y restricciones por entidad}
		\label{tab:crud}
		\centering
		\begin{tabular}{l c c c c}
			\toprule
			\textbf{Entidad}    &
			\textbf{Crear}      &
			\textbf{Leer}       &
			\textbf{Actualizar} &
			\textbf{Eliminar}   \\
			\midrule
			Usuarios    & S\'{i} & S\'{i} & S\'{i}  & S\'{i}  \\
			Productos   & S\'{i} & S\'{i} & S\'{i}  & S\'{i}  \\
			Lotes       & S\'{i} & S\'{i} & S\'{i}  & No* \\
			Movimientos & S\'{i} & S\'{i} & No* & No* \\
			\bottomrule
			\multicolumn{5}{l}{%
				\footnotesize *Restringido para preservar la integridad
				de auditor\'{i}a.}
		\end{tabular}
	\end{table}
	

	%  SECCION 3 — ESTADO DEL ARTE
	
	\section{Estado del Arte}
	\label{sec:arte}
	
	% ------- 3.1 -----------------------------------------------
	\subsection{An\'{a}lisis de Soluciones Existentes}
	\label{subsec:referentes}
	
	Se analizaron tres referentes representativos del mercado:
	
	\begin{itemize}[leftmargin=*, label=\textbullet, itemsep=3pt]
		\item \textbf{Sortly:} aplicaci\'{o}n de gesti\'{o}n visual intuitiva;
		carece de funcionalidades para control de caducidades y
		m\'{e}todo \peps{}. Presenta costo elevado para el segmento
		objetivo.
		
		\item \textbf{Odoo Inventory:} ERP robusto con trazabilidad
		completa, cuya curva de aprendizaje resulta excesiva
		para micronegocios sin personal t\'{e}cnico dedicado.
		
		\item \textbf{Inventory Now:} soluci\'{o}n simple con soporte
		\textit{offline}, pero sin gesti\'{o}n de fechas de
		vencimiento ni soporte para m\'{u}ltiples sucursales.
	\end{itemize}
	
	% ------- 3.2 -----------------------------------------------
	\subsection{Diferenciaci\'{o}n de \app{}}
	\label{subsec:diferenciacion}
	
	\app{} se distingue por su \textbf{enfoque nativo en caducidades y
		en el m\'{e}todo \peps{}}. Integra alertas proactivas dise\~{n}adas
	espec\'{i}ficamente para productos perecederos, combinando simplicidad
	de uso con trazabilidad avanzada a un costo significativamente
	menor que los ERPs disponibles en el mercado.
	
	
	%  SECCION 4 — DEFINICION FUNCIONAL DEL SISTEMA
	
	\section{Definici\'{o}n Funcional del Sistema}
	\label{sec:funcional}
	
	La arquitectura funcional de \app{} est\'{a} dise\~{n}ada para automatizar
	la trazabilidad de productos perecederos. El n\'{u}cleo l\'{o}gico se
	fundamenta en principios cl\'{a}sicos de gesti\'{o}n de inventarios
	aplicados a cadenas de suministro, priorizando la rotaci\'{o}n de
	art\'{i}culos sensibles al tiempo~\cite{Silver:2016}.
	
	% ------- 4.1 -----------------------------------------------
	\subsection{M\'{o}dulos Centrales de Operaci\'{o}n}
	\label{subsec:modulos}
	
	El sistema se divide en cuatro ejes funcionales interconectados:
	
	\begin{itemize}[leftmargin=*, label=\textbullet, itemsep=2pt]
		\item \textbf{Dashboard de control:} panel anal\'{i}tico que consolida
		indicadores clave de rendimiento (KPI) y centraliza alertas
		cr\'{i}ticas en tiempo real.
		
		\item \textbf{Cat\'{a}logo maestro:} m\'{o}dulo de definici\'{o}n de
		art\'{i}culos donde se establecen par\'{a}metros base para
		automatizar alertas de reabastecimiento.
		
		\item \textbf{Gesti\'{o}n de lotes y m\'{e}todo \peps{}:} agrupa las
		entradas en lotes y aplica estrictamente el m\'{e}todo
		\textbf{Primero en Entrar, Primero en
			Salir}~\cite{Silver:2016}, descontando autom\'{a}ticamente
		el lote con fecha de caducidad m\'{a}s pr\'{o}xima en cada
		proceso de salida o venta.
		
		\item \textbf{Motor de trazabilidad:} registro inmutable de cada
		transacci\'{o}n etiquetada con marca de tiempo y autor
		de la acci\'{o}n.
	\end{itemize}
	
	% ------- 4.2 -----------------------------------------------
	\subsection{Flujo L\'{o}gico y Ciclo de Vida del Producto}
	\label{subsec:flujo}
	
	El ciclo de vida de la informaci\'{o}n opera en tres fases,
	ilustradas en la Fig.~\ref{fig:arquitectura}:
	
	\begin{enumerate}[leftmargin=*, label=\arabic*., itemsep=2pt]
		\item \textbf{Registro (entrada):} ingreso de nueva mercanc\'{i}a
		con asignaci\'{o}n obligatoria de fecha de vencimiento.
		
		\item \textbf{Monitoreo activo:} un \textit{cron job} eval\'{u}a
		diariamente cada lote, cambiando su estado de
		\emph{Normal} a \emph{Riesgo} y luego a \emph{Vencido}.
		
		\item \textbf{Ejecuci\'{o}n (salida):} al registrar una venta, la
		l\'{o}gica de base de datos substrae del lote m\'{a}s antiguo,
		garantizando la rotaci\'{o}n \peps{}.
	\end{enumerate}
	
	\begin{figure}[htbp]
		\centering
		\figura{\columnwidth}{4cm}{flujo_datos.png}{Flujo datos PEPS}
		\caption{Flujo l\'{o}gico de informaci\'{o}n y aplicaci\'{o}n del
			m\'{e}todo \peps{}.}
		\label{fig:arquitectura}
	\end{figure}
	
	%  SECCION 5 — FLUJO DE NAVEGACION
	
	\section{Flujo de Navegaci\'{o}n}
	\label{sec:navegacion}
	
	% ------- 5.1 -----------------------------------------------
	\subsection{Mapa Estructural de Vistas}
	\label{subsec:mapa}
	
	El mapa de navegaci\'{o}n garantiza acceso r\'{a}pido a las funciones
	operativas desde cualquier nivel de la jerarqu\'{i}a de pantallas.
	La estructura de vistas se organiza as\'{i}:
	
	\medskip
	\noindent\texttt{LOGIN}\\
	\noindent\texttt{\ \ +-{}- DASHBOARD\ \ (Alertas y m\'{e}tricas)}\\
	\noindent\texttt{\ \ \ \ \ \ +-{}- CATALOGO\ \ \ (B\'{u}squeda / creaci\'{o}n)}\\
	\noindent\texttt{\ \ \ \ \ \ +-{}- LOTES\ \ \ \ \ \ (Entradas y mermas)}\\
	\noindent\texttt{\ \ \ \ \ \ +-{}- PROVEEDORES\ (Contactos y precios)}\\
	\noindent\texttt{\ \ \ \ \ \ +-{}- REPORTES\ \ \ (Auditor\'{i}a)}\\
	\noindent\texttt{\ \ \ \ \ \ +-{}- USUARIOS\ \ \ (Roles y permisos)}
	\medskip
	
	% ------- 5.2 -----------------------------------------------
	\subsection{Identificaci\'{o}n de Acciones del Usuario}
	\label{subsec:acciones}
	
	La Tabla~\ref{tab:acciones} detalla las operaciones cr\'{i}ticas que
	el usuario ejecuta en cada vista para mantener activo el modelo
	\peps{} y el control de caducidades.
	
	\begin{table}[htbp]
		\renewcommand{\arraystretch}{1.35}
		\caption{Acciones operativas y funcionales por vista}
		\label{tab:acciones}
		\centering
		\begin{tabular}{>{\bfseries}p{1.8cm} p{5.6cm}}
			\toprule
			Vista & Acciones clave del usuario \\
			\midrule
			Dashboard
			& Revisar m\'{e}tricas generales y acceder mediante
			atajos a productos pr\'{o}ximos a caducar. \\[3pt]
			Alertas
			& Gestionar notificaciones cr\'{i}ticas y ejecutar
			acciones como aplicar descuentos o registrar
			mermas. \\[3pt]
			Lotes
			& Visualizar el orden de salida \peps{} e ingresar
			mercanc\'{i}a con fecha de vencimiento. \\[3pt]
			Cat\'{a}logo
			& Consultar inventario e identificar visualmente
			el estado del stock (normal, bajo o cr\'{i}tico). \\
			\bottomrule
		\end{tabular}
	\end{table}
	

	%  SECCION 6 — DESARROLLO DE LA INTERFAZ Y PROTOTIPO
	
	\section{Desarrollo de la Interfaz y Prototipo}
	\label{sec:prototipo}
	
	% ------- 6.1 -----------------------------------------------
	\subsection{Dise\~{n}o Centrado en el Usuario y Sem\'{a}ntica Visual}
	\label{subsec:diseno}
	
	Se desarroll\'{o} un prototipo navegable de alta fidelidad basado en
	los lineamientos de \textit{Material Design~3}~\cite{MaterialDesign:2026}.
	El objetivo fue minimizar la carga cognitiva del usuario en entornos
	de alta rotaci\'{o}n comercial. Se implement\'{o} un sistema de
	\textbf{sem\'{a}ntica de colores} estandarizado para guiar decisiones
	operativas en tiempo real:
	
	\begin{itemize}[leftmargin=*, label=\textbullet, itemsep=2pt]
		\item \textbf{Rojo (cr\'{i}tico):} acciones destructivas, productos
		vencidos o quiebre de stock.
		
		\item \textbf{Amarillo / Naranja (advertencia):} lotes que
		entrar\'{a}n en ventana de caducidad en los pr\'{o}ximos 7~d\'{i}as.
		
		\item \textbf{Verde / Neutro (\'{o}ptimo):} estados normales de
		inventario y retroalimentaci\'{o}n positiva ante acciones
		completadas~\cite{MaterialDesign:2026}.
	\end{itemize}
	
	Esta coherencia visual gu\'{i}a la atenci\'{o}n del operario hacia
	las tareas prioritarias, atacando directamente la problem\'{a}tica
	descrita en la Secci\'{o}n~\ref{sec:problema}.
	
	% ------- 6.2 -----------------------------------------------
	\subsection{Vistas Funcionales del Sistema}
	\label{subsec:vistas}
	
	A continuaci\'{o}n se presentan las vistas que conforman el ecosistema
	de la \pwa{} \app{}.
	
	% ··· 6.2.1 -------------------------------------------------
	\subsubsection{Vista General y Alertas Tempranas}
	
	El \textbf{Dashboard} (Fig.~\ref{fig:dashboard}) funciona como
	centro de mando: prioriza las alertas del d\'{i}a y grafica el flujo
	de entradas y salidas, otorgando al administrador visibilidad
	instant\'{a}nea del estado operativo. El \textbf{Centro de Alertas}
	(Fig.~\ref{fig:alertas}) permite ejecutar acciones directas
	desde la notificaci\'{o}n ---como dar de baja un producto o aplicar
	promociones para evitar mermas---.
	
	\begin{figure}[htbp]
		\centering
		\figura{\columnwidth}{5cm}{vista_dashboard.png}{Dashboard principal}
		\caption{Dashboard principal con resumen m\'{e}trico y atajos
			de alerta.}
		\label{fig:dashboard}
	\end{figure}
	
	\begin{figure}[htbp]
		\centering
		\figura{\columnwidth}{5cm}{vista_alertas.png}{Centro de alertas}
		\caption{Centro de alertas clasificado por nivel de criticidad.}
		\label{fig:alertas}
	\end{figure}
	
	% ··· 6.2.2 -------------------------------------------------
	\subsubsection{Control \peps{} y Cat\'{a}logo Visual}
	
	La vista de \textbf{Gesti\'{o}n de Lotes} (Fig.~\ref{fig:lotes})
	organiza la mercanc\'{i}a autom\'{a}ticamente por fecha de caducidad e
	indica qu\'{e} lote debe salir primero seg\'{u}n la prioridad \peps{}.
	El \textbf{Cat\'{a}logo de Productos} (Fig.~\ref{fig:productos}) utiliza
	tarjetas visuales para identificar r\'{a}pidamente la disponibilidad
	y el umbral de stock de cada art\'{i}culo.
	
	\begin{figure}[htbp]
		\centering
		\figura{\columnwidth}{5cm}{vista_lotes.png}{Gestion de lotes PEPS}
		\caption{Gesti\'{o}n de lotes con priorizaci\'{o}n autom\'{a}tica
			del m\'{e}todo \peps{}.}
		\label{fig:lotes}
	\end{figure}
	
	\begin{figure}[htbp]
		\centering
		\figura{\columnwidth}{5cm}{vista_productos.png}{Catalogo de productos}
		\caption{Cat\'{a}logo de productos con indicadores visuales de stock.}
		\label{fig:productos}
	\end{figure}
	
	% ··· 6.2.3 -------------------------------------------------
	\subsubsection{Auditor\'{i}a, Proveedores y Administraci\'{o}n}
	
	El m\'{o}dulo de \textbf{Reportes} (Fig.~\ref{fig:reportes}) documenta
	cada movimiento de forma inmutable, destacando los productos con
	mayores p\'{e}rdidas. El m\'{o}dulo de \textbf{Proveedores}
	(Fig.~\ref{fig:proveedores}) centraliza contactos para facilitar el
	reabastecimiento. La gesti\'{o}n de \textbf{Usuarios}
	(Fig.~\ref{fig:usuarios}) restringe las acciones cr\'{i}ticas al rol
	de administrador.
	
	\begin{figure}[htbp]
		\centering
		\figura{\columnwidth}{5cm}{vista_reportes.png}{Panel de reportes}
		\caption{Panel de reportes con historial de auditor\'{i}a inmutable.}
		\label{fig:reportes}
	\end{figure}
	
	\begin{figure}[htbp]
		\centering
		\figura{\columnwidth}{5cm}{vista_proveedores.png}{Directorio proveedores}
		\caption{Directorio y m\'{e}tricas de rendimiento de proveedores.}
		\label{fig:proveedores}
	\end{figure}
	
	\begin{figure}[htbp]
		\centering
		\figura{\columnwidth}{5cm}{vista_usuarios.png}{Gestion de usuarios}
		\caption{Gesti\'{o}n de usuarios y matriz de permisos por roles.}
		\label{fig:usuarios}
	\end{figure}
	
	\FloatBarrier
	

	%  SECCION 7 — MODELADO DE LA BASE DE DATOS

	\section{Modelado e Implementaci\'{o}n de la Base de Datos}
	\label{sec:bd}
	
	Para dar soporte estructural a la l\'{o}gica de negocio, se dise\~{n}\'{o}
	una base de datos relacional orientada a eliminar redundancias,
	facilitar la escalabilidad y garantizar la integridad de las
	transacciones operativas del comercio.
	
	% ------- 7.1 -----------------------------------------------
	\subsection{Modelo Entidad-Relaci\'{o}n}
	\label{subsec:er}
	
	Antes de revisar las decisiones de dise\~{n}o, es fundamental
	comprender la estructura global sobre la que opera el sistema.
	El Diagrama Entidad-Relaci\'{o}n (Fig.~\ref{fig:er_diagram}) muestra
	c\'{o}mo se relacionan las entidades centrales: productos, lotes,
	movimientos, proveedores y usuarios. Cada caja representa una
	tabla; las l\'{i}neas que las conectan indican que comparten
	informaci\'{o}n, y los s\'{i}mbolos en sus extremos expresan cu\'{a}ntos
	registros de un lado pueden relacionarse con cu\'{a}ntos del otro
	(cardinalidad).
	
	\begin{figure}[htbp]
		\centering
		\figura{\columnwidth}{7cm}{diagramaER_FreshStock.png}{%
			Diagrama ER FreshStock}
		\caption{%
			Diagrama Entidad-Relaci\'{o}n de \app{}. Cada caja es
			una tabla; las l\'{i}neas expresan relaciones y cardinalidades.%
		}
		\label{fig:er_diagram}
	\end{figure}
	
	% ------- 7.2 -----------------------------------------------
	\subsection{Justificaci\'{o}n del Dise\~{n}o}
	\label{subsec:justificacion_bd}
	
	Con la estructura general en mente, se explican a continuaci\'{o}n
	las decisiones de modelado m\'{a}s relevantes.
	
	% ··· 7.2.1 -------------------------------------------------
	\subsubsection{Separaci\'{o}n entre Cat\'{a}logo y Existencias F\'{i}sicas}
	
	La decisi\'{o}n m\'{a}s trascendental fue separar el cat\'{a}logo maestro
	(\texttt{PRODUCTS}) de las existencias f\'{i}sicas (\texttt{BATCHES}).
	Un mismo producto ---por ejemplo, ``Yogurt''--- puede llegar al
	almac\'{e}n en distintas fechas y con distintas fechas de vencimiento.
	Sin esta separaci\'{o}n, ser\'{i}a imposible identificar qu\'{e} unidades
	vencer\'{a}n primero. Al agrupar en lotes, el sistema identifica
	autom\'{a}ticamente cu\'{a}l debe salir antes, ejecutando el m\'{e}todo
	\peps{} con exactitud.
	
	La entidad \texttt{MOVEMENTS} act\'{u}a como libro contable digital:
	registra de forma inmutable cada entrada, salida o merma, junto
	con el usuario responsable y la marca de tiempo exacta.
	
	% ··· 7.2.2 -------------------------------------------------
	\subsubsection{Relaci\'{o}n N:M entre Productos y Proveedores}
	
	Un proveedor suministra m\'{u}ltiples productos; un producto puede
	adquirirse a varios proveedores. Esta relaci\'{o}n de muchos a muchos
	(N:M) se resuelve con la tabla asociativa
	\texttt{PRODUCT\_SUPPLIERS}, que adem\'{a}s almacena el atributo
	\texttt{purchase\_price}: el precio exacto al que cada proveedor
	vende cada producto, informaci\'{o}n esencial para comparar costos
	y calcular m\'{a}rgenes de rentabilidad.
	
	% ··· 7.2.3 -------------------------------------------------
	\subsubsection{Integridad Referencial y Tipos de Datos}
	
	Se configur\'{o} integridad referencial activa: \texttt{ON DELETE
		RESTRICT} impide borrar un producto con lotes asociados, y
	\texttt{ON DELETE CASCADE} elimina registros dependientes cuando
	corresponde. Se emplearon tipos especializados ---\texttt{DECIMAL}
	para montos monetarios y \texttt{TIMESTAMP} para auditor\'{i}a
	con marca de tiempo precisa--- para garantizar exactitud
	en todos los registros operativos.
	

	%  SECCION 8 — ARQUITECTURA DEL BACKEND

	\section{Arquitectura del Backend y L\'{o}gica de Servidor}
	\label{sec:backend}
	
	% ------- 8.1 -----------------------------------------------
	\subsection{Las primeras decisiones del equipo: por qu\'{e} Node.js y Express}
	\label{subsec:decisiones}
	
	Cuando el equipo de desarrollo comenz\'{o} a construir el servidor
	de \app{}, la primera pregunta que se plante\'{o} fue: ``\textquestiondown Qu\'{e}
	tecnolog\'{i}a necesita una fruter\'{i}a real?''. No una
	multinacional, sino un negocio donde tres personas pod\'{i}an
	estar consultando stock, vendiendo y revisando alertas
	exactamente al mismo tiempo. Se dio cuenta de que el servidor
	deb\'{i}a comportarse como un cajero r\'{a}pido: no pod\'{i}a
	dejar la fila en pausa mientras otro cliente decid\'{i}a qu\'{e}
	comprar. Por eso eligi\'{o} \textbf{Node.js}. Su modelo
	\emph{non-blocking I/O} significaba que, si un vendedor ped\'{i}a
	los lotes pr\'{o}ximos a vencer, el servidor no se congelaba
	esperando a la base de datos; atend\'{i}a inmediatamente la
	siguiente petici\'{o}n.
	
	Luego lleg\'{o} \textbf{Express}. Al principio se pens\'{o} en
	escribir todo desde cero, pero r\'{a}pidamente se vio que eso
	implicar\'{i}a perder tiempo parseando cabeceras HTTP y leyendo
	JSON a mano. Express actu\'{o} como un maitre experimentado:
	recib\'{i}a la petici\'{o}n, la clasificaba por URL
	---``/api/products'', ``/api/batches''--- y la derivaba donde
	correspond\'{i}a. Eso permiti\'{o} al equipo concentrarse en la
	l\'{o}gica de negocio en lugar de reinventar la rueda.
	
	% ------- 8.2 -----------------------------------------------
	\subsection{El d\'{i}a que el c\'{o}digo spaghetti oblig\'{o} a reordenar todo}
	\label{subsec:spaghetti}
	
	Al principio del proyecto, todo el backend viv\'{i}a en unos pocos
	archivos que hac\'{i}an de todo. Los DAOs no solo ejecutaban
	consultas SQL, sino que tambi\'{e}n recib\'{i}an \texttt{req} y
	\texttt{res}, formateaban JSON y decid\'{i}an c\'{o}digos de
	estado HTTP. Era como si el cocinero no solo cocinara, sino que
	tambi\'{e}n cobrara en caja y limpiara las mesas. Cuando fue
	necesario corregir un error de stock, el equipo pas\'{o} horas
	perdido entre capas mezcladas. Ah\'{i} entendi\'{o} que el
	principio de responsabilidad \'{u}nica (\textbf{SRP}) no era una
	moda acad\'{e}mica, sino una forma de sobrevivir.
	
	Se decidi\'{o} dividir el backend en tres estaciones bien
	definidas:
	
	\begin{enumerate}[leftmargin=*, label=\arabic*., itemsep=2pt]
		\item \textbf{Routes:} la recepci\'{o}n. Solo anunciaban qu\'{e}
		plato se pod\'{i}a pedir y a qu\'{e} mesa. Por ejemplo:
		\texttt{router.get('/products', productController.getAll)}.
		No sab\'{i}an c\'{o}mo se cocinaba.
		
		\item \textbf{Controllers:} el sal\'{o}n del maitre.
		Recib\'{i}an la orden del cliente, validaban que tuviera
		sentido y llamaban a la cocina. No sab\'{i}an SQL; solo
		sab\'{i}an hablar con el cliente y con los DAOs.
		
		\item \textbf{DAOs:} la l\'{i}nea de cocina. Su \'{u}nica
		preocupaci\'{o}n era ejecutar las recetas SQL en MySQL. No
		sab\'{i}an que exist\'{i}a un protocolo HTTP.
	\end{enumerate}
	
	Esa separaci\'{o}n salv\'{o} al proyecto. Hoy, si hubiera que
	cambiar el motor de base de datos, solo se reescribir\'{i}an los
	DAOs. Los controllers ni siquiera se enterar\'{i}an.
	
	% ------- 8.3 -----------------------------------------------
	\subsection{Los trucos que se aprendieron en el camino: asyncHandler, AppError y el pool}
	\label{subsec:trucos}
	
	Otro dolor de cabeza fue el c\'{o}digo repetido. Cada controller
	ten\'{i}a el mismo bloque \texttt{try/catch}, como un camarero que
	repet\'{i}a ``buenos d\'{i}as, bienvenido'' mec\'{a}nicamente cien
	veces al d\'{i}a. El equipo se cans\'{o} de copiar y pegar, as\'{i}
	que cre\'{o} \texttt{asyncHandler}: un envoltorio de una sola
	l\'{i}nea que atrapaba errores autom\'{a}ticamente y los enviaba
	al middleware central. De pronto, los controllers se le\'{i}an como
	secuencias directas: ``pide los productos, recibe los productos,
	responde los productos''. Sin ruido.
	
	Pero atrapar errores no era suficiente; tambi\'{e}n ten\'{i}an que
	hablar claro. Al principio, cuando un producto no exist\'{i}a, el
	servidor respond\'{i}a un gen\'{e}rico ``Error interno del servidor''
	con c\'{o}digo 500. Eso confund\'{i}a al frontend y volv\'{i}a
	loco al equipo depurando. Se adopt\'{o} entonces una clase
	\textbf{AppError} que portaba su propio c\'{o}digo HTTP: 404 si
	no exist\'{i}a, 400 si los datos eran inv\'{a}lidos, 409 si hab\'{i}a
	duplicados. As\'{i}, cuando un DAO lanzaba \texttt{throw new
	AppError('Producto no encontrado', 404)}, el middleware central
	respond\'{i}a exactamente lo que deb\'{i}a, sin decenas de
	\texttt{if} anidados en cada controller.
	
	Respecto a la base de datos, el equipo aprendi\'{o} que abrir una
	conexi\'{o}n nueva por cada consulta era como si cada mesero fuera
	a la bodega con su propia llave en lugar de usar la puerta
	principal. Cre\'{o} un servicio \texttt{mysql.service.js} que
	inicializaba un \textbf{pool} de conexiones desde el arranque.
	Cuando tres vendedores consultaban stock al mismo tiempo, el
	servidor repart\'{i}a conexiones ya abiertas en lugar de saturar a
	MySQL. Los DAOs simplemente importaban ese pool y ejecutaban
	\texttt{db.query(sql, params)}. Ligero, r\'{a}pido y sin
	conexiones hu\'{e}rfanas.
	
	% ------- 8.4 -----------------------------------------------
	\subsection{Una tarde cualquiera en el servidor: c\'{o}mo fluy\'{o} una petici\'{o}n real}
	\label{subsec:unatarde}
	
	Para que todo esto tuviera sentido, conviene recordar lo que
	ocurri\'{o} una tarde a las tres de la tarde. El encargado del
	almac\'{e}n abri\'{o} su tel\'{e}fono y toc\'{o} ``Ver lotes
	pr\'{o}ximos a vencer''. En ese instante, su navegador envi\'{o}
	una petici\'{o}n \texttt{GET /api/batches/expiring}. Express, en
	el archivo \texttt{batches.routes.js}, la reconoci\'{o} y la
	deriv\'{o} al controller correspondiente.
	
	El controller, envuelto en el \texttt{asyncHandler} que el equipo
	hab\'{i}a creado, extrajo el par\'{a}metro \texttt{days} de la
	URL, valid\'{o} que fuera un n\'{u}mero razonable, y llam\'{o} a
	\texttt{batchDao.getExpiringBatches(days)}. El DAO construy\'{o}
	una consulta SQL con \texttt{JOIN}s sobre las tablas
	\texttt{batch}, \texttt{product} y la funci\'{o}n
	\texttt{DATEDIFF}, y la ejecut\'{o} a trav\'{e}s del pool de
	conexiones.
	
	MySQL devolvi\'{o} las filas en milisegundos. El DAO las entreg\'{o}
	tal cual al controller, quien respondi\'{o}
	\texttt{res.json(rows)} con c\'{o}digo 200. El encargado vio en su
	pantalla los lotes ordenados por fecha de caducidad, exactamente
	como necesitaba. Todo ese recorrido ---desde la URL hasta la base
	de datos y de vuelta--- ocurri\'{o} tan r\'{a}pido que ni siquiera
	perdi\'{o} el ritmo de su recorrido por el almac\'{e}n.
	
	Si algo hubiera fallado ---por ejemplo, que la base de datos
	estuviera ca\'{i}da---, \texttt{asyncHandler} habr\'{i}a atrapado
	la excepci\'{o}n y el middleware de errores habr\'{i}a respondido
	un mensaje claro con c\'{o}digo 500. Sin p\'{a}nico, sin c\'{o}digo
	roto en mil lugares.
	
	El equipo aprendi\'{o} a pensar que cada capa de ese servidor
	dominaba su propio dialecto: HTTP arriba, SQL abajo, y JavaScript
	en el medio como int\'{e}rprete neutral. As\'{i} fue como una
	aplicaci\'{o}n peque\~{n}a pudo crecer sin convertirse en un
	monstruo inmanejable. Eso fue todo lo que hubo detr\'{a}s de
	\app{} cuando nadie lo ve\'{i}a: una arquitectura simple, pero
	pensada para no volverse loca.
	

	%  SECCION 9 — CONCLUSION

	\section{Conclusi\'{o}n}
	\label{sec:conclusion}
	
	% Instruccion \balance para igualar columnas en la ultima pagina
	\balance
	
	En el presente documento t\'{e}cnico se ha definido el alcance
	funcional, visual y de datos de \app{}. La conceptualizaci\'{o}n de
	una \pwa{} respaldada por un dise\~{n}o de interfaces basado en
	sem\'{a}ntica visual valida la propuesta tecnol\'{o}gica para reducir
	mermas en peque\~{n}os comercios de alimentos.
	
	La separaci\'{o}n estrat\'{e}gica de \texttt{PRODUCTS} y \texttt{BATCHES}
	en la base de datos relacional, junto con la normalizaci\'{o}n de
	la relaci\'{o}n con proveedores mediante \texttt{PRODUCT\_SUPPLIERS},
	sientan bases escalables y auditables que dotan al sistema del
	rigor t\'{e}cnico necesario para operar en un entorno real.
	
	De cara a la \'{u}ltima etapa del proyecto, el equipo abordar\'{a} la
	implementaci\'{o}n del sistema de autenticaci\'{o}n segura mediante
	hashing criptogr\'{a}fico propio, el dise\~{n}o e implementaci\'{o}n de la
	capa de persistencia NoSQL bajo un enfoque de \textit{Query Driven
		Design}, y la construcci\'{o}n del proceso ETL que consolidar\'{a} los
	datos operativos en informaci\'{o}n estrat\'{e}gica para la toma de
	decisiones del negocio.

	%  APENDICES

	\appendices
	
	% ------- A -------------------------------------------------
	\section{Recursos del Proyecto}
	\label{apx:recursos}
	
	\begin{itemize}[leftmargin=*, label=\textbullet]
		\item \textbf{Repositorio en GitHub:}\\
		\url{https://github.com/KingSadid/BD_FreshStock.git}
		
		\item \textbf{C\'{o}digo fuente del documento (\LaTeX{}):}\\
		\url{https://www.overleaf.com/4398971276nkbbgsgfvdtb#674337}
	\end{itemize}
	
	% ------- B -------------------------------------------------
	\section{Glosario de T\'{e}rminos Clave}
	\label{apx:glosario}
	
	\begin{itemize}[leftmargin=*, label=\textbullet, itemsep=2pt]
		\item \textbf{PWA (\textit{Progressive Web App}):} aplicaci\'{o}n
		web que ofrece experiencia equivalente a app nativa sin
		instalaci\'{o}n desde tienda de aplicaciones.
		
		\item \textbf{PEPS:} ``Primero en Entrar, Primero en Salir''
		(\textit{FIFO}). M\'{e}todo de rotaci\'{o}n de inventarios
		para productos perecederos.
		
		\item \textbf{Merma:} p\'{e}rdida f\'{i}sica de productos por
		caducidad, da\~{n}o o deterioro.
		
		\item \textbf{Modelo ER:} diagrama que representa la estructura
		conceptual de una base de datos mediante entidades,
		atributos y relaciones.
		
		\item \textbf{KPI (\textit{Key Performance Indicator}):}
		indicador clave de rendimiento operativo.
		
		\item \textbf{CRUD:} operaciones b\'{a}sicas sobre datos:
		\textit{Create}, \textit{Read}, \textit{Update},
		\textit{Delete}.
		
		\item \textbf{ETL (Extract, Transform, Load):} proceso que
		extrae datos operativos, los transforma en informaci\'{o}n
		consolidada y los carga en estructuras optimizadas
		para anal\'{i}tica.
		
		\item \textbf{NoSQL:} bases de datos no relacionales, optimizadas
		para grandes vol\'{u}menes y estructuras de datos flexibles.
		
		\item \textbf{Query Driven Design:} metodolog\'{i}a de modelado
		NoSQL que dise\~{n}a la estructura a partir de las
		consultas m\'{a}s frecuentes del sistema.
		
		\item \textbf{Hashing criptogr\'{a}fico:} transformaci\'{o}n
		irreversible de una contrase\~{n}a mediante una funci\'{o}n
		matem\'{a}tica, de modo que la contrase\~{n}a original
		nunca se almacene en texto plano.
	\end{itemize}
	

	%  REFERENCIAS

	\begin{thebibliography}{1}
		
		\bibitem{Silver:2016}
		E.~A. Silver, D.~F. Pyke, and D.~J. Thomas,
		\emph{Inventory and Production Management in Supply Chains},
		4th~ed.
		Boca Raton, FL, USA: CRC Press, 2016.
		
		\bibitem{PMI:2019}
		Project Management Institute,
		``Small Business Inventory Management: Best Practices,''
		PMI White Paper, 2019.
		
		\bibitem{FAO:2019}
		Food and Agriculture Organization of the United Nations,
		\emph{The State of Food and Agriculture 2019: Moving Forward on
			Food Loss and Waste Reduction},
		Rome, Italy: FAO, 2019.
		
		\bibitem{MaterialDesign:2026}
		Google LLC,
		``Material Design 3,'' 2026.
		[Online]. Disponible en: \url{https://m3.material.io/}
		
	\end{thebibliography}
	
\end{document}